\section{力学相对性原理与伽利略变换}
\subsection{参考系与坐标系}
\textbf{参考系}是描述物体运动时被选作参照的物体或物体系，它决定了运动的观测结果。
\textbf{坐标系}是固定在参考系上的数学工具，用于定量描述位置和方向。
同一参考系可采用不同的坐标系.参考系是物理的，而坐标系是数学的。

例如，当我们坐在行驶的汽车上观察路边的树时，我们以汽车为参考系，就会看到树以速度$\mathbf v$向我们的后方移动，无论采取何种坐标系，都不能改变树木向我们的后面移动的这种现象。但是，我们可以采用不同的坐标系来描述这个速度矢量$\mathbf v$，如在正方向向前的一维坐标系内，则$\mathbf v = - v \mathbf i$，这里的$\mathbf i$为向前方的单位向量；当然也可以用正方向向后的一维坐标系，那么$\mathbf v = v \mathbf i '$，这里的$\mathbf{ i } '$为向后方的单位向量。我们通过取不同的坐标系就可以让同样的矢量具备不同的分量。

不过坐标系可以很好地帮助我们传递信息。例如我们描述某人正在以速度$v$向西运动，就依靠了东南西北上下这样一个坐标系，这比起一个不依赖于坐标系的矢量$\mathbf v$更具有现实意义。

总结来说，参考系和坐标系是截然不同的两个概念，需要我们严格辨析。我们许多的物理定律是可以不依赖于坐标系的，想要准确描述其规律却可能需要在一个坐标系中才能比较清晰。
\subsection{伽利略变换}
点P在空间中运动，我们尝试在两个惯性系$K$和$K'$中考察其运动。设惯性系 \( K' \) 相对惯性系 \( K \) 以恒定速度 \( u \) 运动，两参考系原点分别为$O'$和$O$。\( P \) 点在 \( K \) 系中的位矢是 \( \mathbf{r} \)，在 \( K' \) 系中的位矢是 \( \mathbf{r}' \)。为简单起见，假定 \( t = 0 \) 时刻两参考系原点重合，则 \( \overrightarrow{OO'} = \mathbf{u} t \)，从而得到如下关系式

\[
\begin{cases} 
\mathbf{r} = \mathbf{r}' + \mathbf{u} t, \\ 
\mathbf{v} = \mathbf{v}' + \mathbf{u}, \\ 
\mathbf{a} = \mathbf{a}'.
\end{cases}
\]
这一式子的形式不依赖于坐标轴的选取。

如果我们再分别在两个参考系中选取两个空间直角坐标系，两者的原点分别固定于$K$和$K'$，$x$轴正方向与$\mathbf u$平行，$y$与$y'$、$z$与$z'$平行，那么上面的变换关系就可以写为如下坐标变换
\[\begin{cases} 
x = x' + ut, \\ 
y = y', \\ 
z = z', \\ 
t = t'. 
\end{cases}\]
这组变换被称为\textbf{伽利略变换}。
\subsection{力学相对性原理}
在经典力学中，如果两个参考系之间没有相对的加速度，且已知其中之一为惯性系，那么另一个也必然是惯性系。这两个惯性系之间是完全对称的，没有谁优谁劣的差异。伽利略提出的\textbf{力学相对性原理}，表述为
\begin{center}
力学规律对一切惯性系等价
\end{center}

我们可以这么理解这一原理。由于我们考虑的两个参考系均为惯性系，彼此之间相对速度恒定，根据上面的伽利略变换，一个物体在两个参考系中的加速度将会一致。这样，如果要求牛顿第二定律$F=ma$成立，两参考系内必然同时具有完全相同的力学规律。

这一原理在数学逻辑上是完全自洽的，但是如果我们直接把该原理表述为“物理规律对一切惯性系等价”，就会在电磁学领域出现一定问题。为了能够涵盖电磁学，保持物理规律的简洁性，我们不能改动动力学规律，所以我们不得不修改我们对于坐标与时间的认识，认为满足相对性原理的坐标变换为洛伦兹变换而非伽利略变换，这就是\textbf{绝对时空观}和\textbf{相对论时空观}的差异。